\documentclass[10pt,a4paper]{amsart}
\usepackage[margin=19mm]{geometry}
\usepackage{amsmath,amssymb,mathtools}
\usepackage[scr=rsfs,cal=euler]{mathalfa}
\usepackage{xfrac}
\usepackage[unicode,hidelinks,bookmarks=false]{hyperref}
\newcommand{\ie}{i.e.\ }
\newcommand{\cf}{cf.\ }
\newcommand{\resp}{respectively\ }
\newcommand{\Subsection}{\subsection}
\providecommand{\nobreakdash}{\nobreak\hbox{-}}
\setlength{\emergencystretch}{2em}
\multlinegap0pt
\usepackage{bm}
\usepackage{bbm}
\usepackage{dsfont}
\usepackage{xspace}
\usepackage{cancel}
\usepackage{mathtools}
\usepackage{amsthm}
\makeatletter
\@ifundefined{lemma}{\newtheorem{lemma}{Lemma}}{}
\@ifundefined{prop}{\newtheorem{prop}{Proposition}}{}
\makeatother
\title[Infinite systems of equations in abelian and nilpotent group]{Infinite systems of equations in abelian and nilpotent groups}
\author{Mikhail A. Mikheenko}
\date{Source: arXiv:2410.20729v3 (CC BY 4.0)}
\begin{document}
\maketitle

\noindent\emph{Source and licence.} Infinite systems of equations in abelian and nilpotent groups. Version arXiv:2410.20729v3 (CC BY 4.0), distributed under the Creative Commons Attribution 4.0 International licence, as recorded in the arXiv licence metadata for this exact version and shown on the abstract page https://arxiv.org/abs/2410.20729v3. Full rights notice: licenses/arxiv-2410.20729-CC-BY-4.0.txt.\par\smallskip

\noindent\textit{Editorial scope.} Counted unit: the Lemma whose source label is \texttt{bad} in the article (source0.tex lines 719--725, source label \texttt{bad}), delivered together with the article's own complete original proof of it (source0.tex lines 726--780, one \texttt{proof} environment of 55 lines). No mathematics is written, repaired or re-derived: statement and proof are the article's own text line for line. Required context retained uncounted: pbad (lines 602--607). Every definition, display and label of those lines is retained unchanged. Omitted from this package and not redistributed: 1--601, 608--718, 781--1457. Nothing inside a retained excerpt is omitted or altered: every retained body is delivered exactly as the article's lines are written, so no omission receipt of any kind accompanies this package. Citations: the retained excerpts carry no citation command at all, so no bibliography entry of the article is transcribed and no reference is redistributed. Reference adaptation: none was needed and none was made; every reference inside the delivered unit resolves inside it, so no unresolved reference or citation is produced. Printed slips kept literal: no printed word, symbol, hypothesis or number of the counted statement or of its proof was altered. Layout and class: the article's own class is \texttt{article} with packages including cmap, amsthm, inputenc, fontenc, amsfonts, amssymb, amsmath, color, hyperref, mathdots, cite, geometry; the wrapper is the lane's pinned \texttt{amsart} editorial wrapper at 10pt on A4, which restates the theorem-like environments the retained text needs and adds the article's own macro definitions that the retained text needs (none), with extra wrapper packages (bm, bbm, dsfont, xspace, cancel, mathtools). The delivered numbering is the unit's own. Assets: no retained span contains a figure environment, a graphics-inclusion command, a PDF-page inclusion, a table or a TikZ picture; no image file is copied into this package, the package declares no asset dependency and no image file is redistributed with it. Licence: version arXiv:2410.20729v3 of this article is distributed under the Creative Commons Attribution 4.0 International licence, as recorded in the arXiv licence metadata for this exact version and shown on the abstract page https://arxiv.org/abs/2410.20729v3. A full rights notice is delivered at licenses/arxiv-2410.20729-CC-BY-4.0.txt. Characters: every retained excerpt is pure ASCII, so the delivered engine, XeLaTeX with its default Latin Modern text font, reports no missing character.
\begin{prop}\label{pbad}
Suppose that $A$ is a reduced abelian $p$-group.
Suppose that the period of $A$ is not bounded.
Then there is a unimodular system of equations over $A$,
which has no solutions in $A$.
\end{prop}

\begin{lemma}\label{bad}
Suppose that $A$ is a reduced periodic abelian group
of unbounded period.
Then there is a unimodular system of equations
over $A$
which has no solutions in $A$.
\end{lemma}

\begin{proof}
The periodic abelian group $A$ is a direct sum
of its $p$-components:
$$
A = \bigoplus_{i \in I} A_{p_i}.
$$
First consider the case where $I$ is finite.
Then some component $A_p$ has unbounded period
(otherwise the period of $A$ is bounded).
If any unimodular system of equations over $A$ has a solution in $A$,
then this is true for any quotient group of $A$.
However, Proposition \ref{pbad} shows that
for $A_p$ this fact is not true.
So, not all unimodular systems of equations over $A$
are solvable in $A$. One can even give an example
of a system which has no solutions in $A$: if the projection
of the system onto $A_p$ is the system from the proof of Proposition \ref{pbad},
then the system itself is unsolvable in $A$.

Now suppose that $I$ is infinite (in this case we can assume $I = \mathbb N$,
although the numbers $p_i$ do not have to range over all prime numbers).
In each $A_{p_i}$ pick an element
$a_i \in A_{p_i}\setminus p_iA_{p_i}$, which exists
since the abelian $p_i$-group $A_{p_i}$ is not divisible.
Consider the following system of equations:
$$
\{x + p_i y_i = a_i \mid i \in \mathbb N \}.
$$
It has the following matrix of exponent sums:
$$
\begin{pmatrix}
1 & p_1 & 0 & 0 &\dots \\
1 & 0 & p_2 & 0 &\dots \\
1 & 0 & 0 & p_3 & \dots \\
\dots & \dots & \dots & \dots & \dots \\
\end{pmatrix}
$$
Clearly, the system is $p_i$-nonsingular
for every $p_i$ and $p$-nonsingular
for the remaining prime numbers $p$, so this system is unimodular.

Suppose that this system has a solution $\{\tilde x, \tilde y_1, \tilde y_2, \ldots\}$
in $A$. It means that this solution lies in the Cartesian product
$\prod\limits_{i \in \mathbb N} A_{p_i}$ of $p$-components.
Look at the component $A_{p_i}$ and the equation $\tilde x + p_i \tilde y_i = a_i$.
As $a_i \notin p_i A_{p_i}$, $p_i$-component of $\tilde x$
is not $0$ (otherwise $a_i$ is equal to $p_i \hat y_i \in p_iA_{p_i}$, where
$\hat y_i$ is the $p_i$-component of $\tilde y_i$).
So we get that every component of $\tilde x\in\prod\limits_{i \in \mathbb N} A_{p_i}$
is not equal to $0$.
As there is an infinite number of such components in $\prod\limits_{i \in \mathbb N} A_{p_i}$,
$\tilde x$ does not lie in the direct sum of the components, which means $\tilde x \notin A$.
We get a contradiction, so the assumption of finding a solution
of the system in $A$ is false.
\end{proof}

\end{document}
